\preliminar{Resumen}

El presente proyecto aborda la ineficiencia del proceso de venta y reserva
de paquetes turísticos en Vizcacha Tours, Potosí, donde la verificación de
cupos y la actualización de la información no se realizan de manera uniforme
entre el personal. El objetivo general es desarrollar un sistema web de
gestión de reservas que mejore la consistencia del proceso de venta y reserva
de paquetes turísticos en Vizcacha Tours. La propuesta se sustenta en los
antecedentes de los sistemas de reservas, en el marco de la gestión de
paquetes turísticos y en criterios de usabilidad y seguridad para
aplicaciones web. La metodología prevé entrevistas, observación y revisión
autorizada de documentos en la agencia, con el fin de diagnosticar el proceso
actual y diseñar la solución de acuerdo con sus necesidades. El sistema será
utilizado por el personal de la agencia y por los clientes, que accederán con
una cuenta propia para consultar la disponibilidad, solicitar reservas y
registrar sus pagos; ambas entradas siguen el mismo criterio de validación y
de confirmación. Se espera que el sistema contribuya a reducir inconsistencias
en la disponibilidad, disminuir el tiempo de confirmación y ofrecer
información confiable para la toma de decisiones.

\vspace{0.5\baselineskip}
\noindent\textbf{Palabras clave:} agencias de viaje, paquetes turísticos,
reservas, sistemas web, control de disponibilidad
